% Testo della pagina divisoria "Fuoco di bivacco" (inserito da generate_latex.py).
\begin{center}\begin{minipage}{0.86\textwidth}
Un buon fuoco ha un andamento a onda: sale, tocca il punto più alto e poi scende fino al silenzio.
I canti sono raccolti per fase, nell'ordine in cui si cantano.

\medskip
\begin{tabular}{@{}l@{\hspace{4mm}}p{0.72\textwidth}@{}}
\textbf{Apertura} & Si accende il fuoco: i canti che lo introducono.\\[1mm]
\textbf{Bans} & Brevi, con gesti: sciolgono il gruppo e danno ritmo.\\[1mm]
\textbf{Crescendo} & Canti noti e ritmati, la voce sale.\\[1mm]
\textbf{Picco} & Il momento più acceso, breve.\\[1mm]
\textbf{Calando} & Ritmo e voce scendono: canti di strada e di ricordo.\\[1mm]
\textbf{Chiusura} & Canti raccolti, quasi sottovoce.\\[1mm]
\textbf{Ultimi} & \textit{Scende la notte} e \textit{Signor, fra le tende schierati}, poi silenzio.
\end{tabular}

\medskip
\textit{Regola d'oro:} chi guida legge il fuoco, non la scaletta. Se il gruppo è stanco, si scende prima.
\end{minipage}\end{center}
